\documentclass[../../main.tex]{subfiles}
\begin{document}
\section{Family 2: stochastic localization}
\label{sec:family-sl}

\paragraph{Object followed.}
A measure-valued martingale $t\mapsto\mu_t$, together with its centroid $a_t$ and covariance
$A_t$, along a Brownian filtration.

\paragraph{What it buys.}
Every general improvement of the KLS bound since 2012 has come from this construction. It replaces
the needle dichotomy of Section~\ref{sec:family-needles} --- ``one dimension, no covariance'' ---
by a process that keeps the ambient dimension and tracks covariance explicitly, at the price of
controlling it only in law.

This section fixes the process and the three identities used by the stochastic routes in
Groups~3--4. The
repository's own two-color refinement of them is Section~\ref{sec:notation} onward; what follows
is the scalar/matrix backbone as it appears in the literature.

\subsection{The process}
\label{subsec:sl-process}

In the Lee--Vempala form of Eldan's process \cite{LeeVempala2018,LeeVempala2024}, start from
$p_0=p$ and define
\begin{equation}
  p_t(x)=\frac{\exp\bigl(\inner{c_t}{x}-\frac t2|x|^2\bigr)\,p(x)}
              {\int\exp\bigl(\inner{c_t}{y}-\frac t2|y|^2\bigr)\,p(y)\dd y},
  \label{eq:sl-density}
\end{equation}
where
\[
  \dd c_t=\dd W_t+a_t\dd t,
  \qquad
  a_t=\E_{p_t}X,
  \qquad
  A_t=\Cov(p_t).
\]

It\^o calculus gives the three identities that carry the whole method:
\begin{align}
  \dd p_t(x)&=p_t(x)\inner{x-a_t}{\dd W_t},
  \label{eq:sl-density-sde}\\[2pt]
  \dd a_t&=A_t\dd W_t,
  \label{eq:sl-centroid-sde}\\[2pt]
  \dd A_t&=\calT_t(\dd W_t)-A_t^2\dd t,
  \label{eq:sl-covariance-sde}
\end{align}
where the matrix-valued third-moment tensor is
\begin{equation}
  \calT_t(u)=\E_{p_t}\bigl[(X-a_t)^{\otimes2}\inner{X-a_t}{u}\bigr].
  \label{eq:third-moment-tensor}
\end{equation}

Two features are decisive, and they pull in opposite directions:

\begin{enumerate}[label=(\roman*),leftmargin=2.4em]
  \item By \eqref{eq:sl-density-sde}, $p_t(x)$ is a \emph{measure-valued martingale}:
  $\E p_t(x)=p_0(x)$ for every $x$. Whatever is proved about $p_t$ can therefore be averaged back
  to time zero.
  \item The factor $e^{-t|x|^2/2}$ in \eqref{eq:sl-density} makes $p_t$ \emph{$t$-strongly
  log-concave}. Curvature is manufactured, at a known rate, out of nothing.
\end{enumerate}

Feature (ii) creates the good event; feature (i) transports it back. The entire difficulty is that
the drift $-A_t^2\dd t$ in \eqref{eq:sl-covariance-sde} that would keep $A_t$ small competes with
the noise $\calT_t(\dd W_t)$, which is a third-moment quantity.

\begin{remark}[Filtering interpretation]
\label{rem:sl-filtering}
The process is a nonlinear filter. The observation $c_t$ has the law of $tX+W_t$, and $p_t$ is
exactly the posterior law of $X$ given that noisy observation up to time $t$. Localization is
therefore ``observe $X$ through Gaussian noise of decreasing variance $1/t$ and watch the
posterior sharpen''. Under this reading $A_t$ is a conditional covariance, which is what makes
Proposition~\ref{prop:covariance-spike} directly relevant: it is a statement about
$\Cov(X\mid X+\sqrt sG)$.
\end{remark}

\subsection{How isoperimetry is transferred back}
\label{subsec:sl-transfer}

Fix a Borel set $E$ and let $g_t=p_t(E)$. By \eqref{eq:sl-density-sde}, $g_t$ is a martingale with
\begin{equation}
  \dd g_t=\Bigl\langle\int_E(x-a_t)\,p_t(x)\dd x,\ \dd W_t\Bigr\rangle,
  \label{eq:mass-martingale-sde}
\end{equation}
whose quadratic variation is controlled by $A_t$. Lee--Vempala's criterion is the clean statement
of the transfer \cite{LeeVempala2018}:
\begin{equation}
  \Prob\Bigl\{\int_0^T\norm{A_s}_\op\dd s\le\tfrac1{64}\Bigr\}\ge\tfrac34
  \qquad\Longrightarrow\qquad
  \PsiKLS_{p_0}\lesssim T^{-1/2}.
  \label{eq:lv-criterion}
\end{equation}

The mechanism has three steps, and it is worth separating them because the routes of Groups~3--4
modify exactly one of them each:
\begin{enumerate}[label=(\arabic*),leftmargin=2.4em]
  \item the covariance bound keeps $g_T$ away from $0$ and $1$ with positive probability --- the
  cut is not identified too fast;
  \item $p_T$ is $T$-strongly log-concave, hence $h(p_T)\gtrsim\sqrt T$ by Bakry--\'Emery;
  \item averaging the boundary measure of $E$ under $p_T$ back to time zero, using the martingale
  property, transfers that expansion to $p_0$.
\end{enumerate}

So KLS becomes a question about the covariance process --- and specifically about its largest
eigenvalue, which is where \eqref{eq:lv-criterion} charges its cost. Route~E replaces step~(1) by
a two-color estimate for one fixed $E$ (Section~\ref{sec:mass-martingale}); Route~S replaces the
set $E$ by a first eigenfunction (Section~\ref{sec:spectral-route}).

\subsection{Why third moments appear}
\label{subsec:sl-third-moments}

Take the basic Schatten-2 potential $\Phi_t=\Tr(A_t^2)$. It\^o's formula applied to
\eqref{eq:sl-covariance-sde} gives
\begin{equation}
  \dd\Phi_t
  =\dd M_t
  -2\Tr(A_t^3)\dd t
  +\E_{X,Y\sim p_t}\inner{X-a_t}{Y-a_t}^3\dd t ,
  \label{eq:trace-potential-ito}
\end{equation}
with $M_t$ a martingale. The term $-2\Tr(A_t^3)$ is the stabilizing drift inherited from $-A_t^2$.
The positive It\^o correction is exactly the squared Hilbert--Schmidt norm of the third-moment
tensor: writing $a=a_t$,
\begin{equation}
  \E_{X,Y\sim p_t}\inner{X-a}{Y-a}^3
  =\sum_{i,j,k}\Bigl(\E\bigl[(X-a)_i(X-a)_j(X-a)_k\bigr]\Bigr)^2 .
  \label{eq:third-moment-identity}
\end{equation}

Controlling this noise term \emph{while retaining information about $\lmax(A_t)$} is the central
stochastic-localization calculation, and it is why a sharp bound on the third-moment tensor ---
the parameter $\kappa_n$ of \eqref{eq:kappa-def} --- is the crucial input rather than an
incidental one.

\subsection{Evolution of the method}
\label{subsec:sl-evolution}

\begin{center}
\begin{tabularx}{\textwidth}{@{}>{\raggedright\arraybackslash}p{0.20\textwidth}Y>{\raggedright\arraybackslash}p{0.24\textwidth}@{}}
\toprule
Paper & Main technical device & Result for $\PsiKLS_n$ \\
\midrule
Eldan \cite{Eldan2013ThinShell} & Covariance-normalized localization $C_t=A_t^{-1}$, high Schatten
powers, relation to thin shell & $n^{1/3}\sqrt{\log n}$ with then-known thin-shell bounds \\
\addlinespace
Lee--Vempala \cite{LeeVempala2018,LeeVempala2024} & Fixed Gaussian tilt and the $\Tr(A_t^2)$
stopping argument & $n^{1/4}$ \\
\addlinespace
Chen \cite{Chen2021} & Iterated high-Schatten potentials, affine preconditioning, induction on
dimension & $\exp(C\sqrt{\log n\log\log n})$ \\
\addlinespace
Klartag--Lehec \cite{KlartagLehec2022Polylog} & Heat-flow duality, spectral projections, $H^{-1}$,
covariance growth & $(\log n)^5$ \\
\addlinespace
Jambulapati--Lee--Vempala \cite{JambulapatiLeeVempala2022KLS} & Two localization representations
and spiked-spectrum estimates & $(\log n)^{3.2226}$ \\
\addlinespace
Klartag \cite{Klartag2023Logarithmic} & Improved Lichnerowicz plus short-time covariance control &
$\sqrt{\log n}$ (published record) \\
\addlinespace
Letwin, v1 \cite{Letwin2026QuadraticKLS} & Sharp quadratic Poincar\'e $\Rightarrow\kappa_n=O(1)$,
inserted into Klartag's bridge & $(\log n)^{1/4}$ (preprint record) \\
\bottomrule
\end{tabularx}
\end{center}

\subsection{Where the surviving logarithm lives}
\label{subsec:sl-where-the-log-lives}

This subsection makes precise the claim of Section~\ref{subsec:kls-architecture}, because it is
what target~3 of the synthesis acts on.

To use \eqref{eq:lv-criterion} one needs a potential that both dominates $\lmax(A_t)$ and admits
an It\^o calculus. Klartag--Lehec use log-trace-exp,
\begin{equation}
  \Phi_t=\frac1\beta\log\Tr\bigl(e^{\beta A_t}\bigr),
  \qquad
  \norm{A_t}_\op\le\Phi_t\le\norm{A_t}_\op+\frac{\log n}\beta .
  \label{eq:logtraceexp}
\end{equation}
The two-sided bound in \eqref{eq:logtraceexp} is the whole story. Approximating $\lmax$ to within
a constant requires $\beta\asymp\log n$. The It\^o drift generated by the noise
\eqref{eq:third-moment-tensor} then has size $O(\kappa_n^2\log n)$, so covariance control survives
only up to
\begin{equation}
  t_*\asymp\frac1{\kappa_n^2\log n} .
  \label{eq:tstar}
\end{equation}
Improved Lichnerowicz (Theorem~\ref{thm:improved-lichnerowicz}) converts a curvature time into a
Poincar\'e constant,
\[
  \CP(\mu)\lesssim t_*^{-1/2}\lesssim\kappa_n\sqrt{\log n},
\]
which is the bridge \eqref{eq:kls-bridge}.

\begin{remark}[The logarithm is an entropy cost, not a moment cost]
\label{rem:log-is-entropy}
The $\log n$ in \eqref{eq:logtraceexp} is the entropy of the uniform distribution on $n$
directions: it is what one pays to replace a maximum over $n$ eigenvalues by a smooth surrogate,
and it would be present even if $\kappa_n$ were known exactly and equal to $1$. Since
$\kappa_n=O(1)$ is now available (Proposition~\ref{prop:letwin-kappa}), this term is the \emph{entire}
remaining gap between the preprint record and the conjecture. A potential that depends only on the
directions actually relevant to a near-extremizer, or on an effective rank rather than the ambient
dimension $n$, would convert $\kappa_n=O(1)$ into $\CP=O(1)$.
\end{remark}

\paragraph{The precise missing estimate.}
Control of $\lmax(A_t)$ along the whole path at a universal time, without paying the $\log n$
of \eqref{eq:logtraceexp}.

\paragraph{Why it stalls.}
By Proposition~\ref{prop:covariance-spike}, the naive strengthening --- a uniform pathwise bound
on $\norm{A_t}_\op$ --- is false, even for measures that satisfy KLS. So the missing estimate
cannot be obtained by sharpening the covariance bound; it has to come from a potential that
recognizes when a covariance spike is harmless. Products of centered exponentials are the
canonical instance of a harmless spike, which is why they recur as the stress test throughout
Groups~3--4 (Sections~\ref{sec:models} and~\ref{sec:product-stress}).
\end{document}
